\section{绪论}
\label{sec:introduction}

印制电路板（Printed Circuit Board，PCB）承载电子器件，并负责电气连接与信号传输。线路开路、短路、鼠咬、毛刺、针孔和多余铜等表面缺陷会改变局部导通状态。这些缺陷还会在后续装配与服役阶段引发间歇性故障。线路密度和制造节拍提高后，人工逐板检查难以兼顾一致性、速度与成本。自动光学检测（Automated Optical Inspection，AOI）已成为 PCB 制造质量控制的常用环节。其典型流程先快速定位疑似区域，再由规则、模型或操作员复判报警是否成立\cite{deng2018autovrs,chen2023pcbreview}。误报会增加停线和人工确认成本，漏检会使结构性缺陷进入后续工序。复核系统因此需要同时评价纠错收益、正确结果保护与处理开销。

基于卷积神经网络和视觉 Transformer 的目标检测器提高了 PCB 缺陷定位的自动化程度\cite{tang2019deeppcb,chen2023pcbreview}。这类方法通常以待检图像为主要输入，通过多尺度特征识别局部纹理异常。单图检测仍有结构性局限。PCB 缺陷往往只占几十个像素，细小目标容易在下采样过程中丢失。正常线路的局部纹理重复、密集且高度相似，灰尘、反光和轻微成像偏差也会产生类似缺陷的响应。开路与短路的判别依据并非普通物体的外观语义，而是线路是否保持应有的连通、隔离和分支关系。仅凭当前图像，系统难以稳定区分多余铜、正常焊盘与光照扰动。叠加检测头、注意力模块或损失项可以改善特征表达，但无法补充同位置的正常状态。

正常模板为这些问题提供了反事实参照。单图检测器回答“该区域是否像缺陷”，模板比较则回答“该位置相对正常设计发生了什么变化”。对于开路，正常模板给出本应连续的铜线；对于短路，正常模板给出本应隔离的线路。对于鼠咬、毛刺、针孔和多余铜，正常模板分别给出完整边缘、正常轮廓、孔结构及铜区范围。配对正常模板使变化判断受到正常结构约束。已有 PCB 检测和工业变化检测工作表明模板信息具有应用价值\cite{tang2019deeppcb,fridman2021changechip}。模板差分仍会受到配准误差、光度变化和阈值噪声影响，因此差异区域只用于提出待验证假设，不能直接视为真实缺陷。

本文将多模态大模型（Multimodal Large Language Model，MLLM）简称为 MLLM。MLLM 能够联合处理局部图像、正常模板和封闭缺陷定义\cite{jiang2025mmad,xu2025anomalyov,andrade2026agreementbias}。与固定阈值规则相比，MLLM 可以组合多种视觉线索，并在预先限定的类别空间内输出结构化判断。通用 MLLM 对极小异常的感知能力有限，也会依赖文字先验或生成不存在的结构。其自报置信度不能直接作为工业决策依据。让 MLLM 在整幅 PCB 图像上自由定位会增加调用成本，并引入不受约束的新框。把基础检测器的类别和置信度直接写入提示，则会产生既有答案锚定风险。因此，本文把 MLLM 限定为局部、选择性、可审计且受安全规则约束的多模态验证器。

本文研究挂接在基础检测器之后的 PCB 缺陷结果复核。系统将检测候选与正常模板差异模块产生的差异候选合并为统一候选空间。多模态验证器只判断已有候选，不在任意位置生成新框。复核调用选择门（简称复检门）从中选择需要调用 MLLM 的高风险候选，接受门再判断提议动作是否有足够的外部证据。前者控制调用效率，后者约束修改风险。最终系统使用 A2 三面板证据卡，包括当前紧裁剪、正常模板紧裁剪和更大上下文。外观差异图与拓扑差异图不作为 MLLM 的视觉输入，拓扑向量仅用于门控决策。系统设计和实验评价均明确考察 MLLM 的能力边界。

本文的贡献包括以下三点。

\begin{enumerate}
  \item \textbf{面向漏检恢复的双来源候选机制。}该机制联合检测候选与差异候选，在统一候选空间中处理误报、错分类和检测器未提出候选的漏检。差异候选的漏检恢复候选率为 0.625（IoU 0.3）。候选上界 F1 为 0.9555，相对基线 0.9118 留有 $+0.044$ 空间，表明候选覆盖并非主要限制因素。五折错误分解中，漏检占可纠错空间的 57.6\%，在三类错误中占比最高。

  \item \textbf{风险受控的分级复核机制。}学习式复检门控制调用预算，外部证据驱动的接受门约束修改动作。复检门在 5\% 调用预算下定位 46.5\% 的可纠错错误，随机策略为 3.2\%，前者是后者的 14.8 倍。有接受门的预算曲线 F1 为 0.9115 至 0.9117，与基线 0.9118 无统计差异；无接受门时，F1 随预算从 0.9088 降至 0.8660。双门控的正确检测保留率为 0.9998，734 个正确检测中误删 2 个。双门控还将有害修正率从无门控时的 0.9326 降至 0.4000，并把误删控制在 2 次。

  \item \textbf{MLLM 作为工业细粒度缺陷验证器的能力边界实证。}正常模板是证据消融中唯一稳定的增益来源。A2 相对 A0 提升 0.1178，置信区间不跨 0。阶段 C 的 200 候选主消融中，A4 的 F1 为 0.7571，略高于 A2 的 0.7553，但该差异在置信区间内，未证明拓扑图带来增量。阶段 B2 的 450 次专门对照中，A2 的 F1 为 0.8082，高于全部含差异图的变体。对渲染自动拉伸与配准残差的排查没有改变这一结论。Oracle 分析显示，候选上界为 0.9555，接受上界为 0.9143，仅比基线高 0.0025。主要限制来自 MLLM 提议动作中正确项不足。附加对照还显示，去锚定未见收益，配准不可靠提示产生负效应，自由文本输出的 JSON 可解析率为 0。
\end{enumerate}

第~\ref{sec:related} 节梳理 PCB 缺陷检测、正常模板比较、工业异常多模态理解及选择性决策研究。第~\ref{sec:problem} 节定义双来源候选、六种复核动作及整体数据流。第~\ref{sec:candidates} 节介绍模板处理、差异候选生成、候选融合和三面板证据卡。第~\ref{sec:gating} 节说明复检门、多模态验证器、动作映射、接受门和最终融合。第~\ref{sec:experiments} 节报告实验设置、证据消融、端到端纠错、预算曲线与鲁棒性分析。第~\ref{sec:conclusion} 节总结研究结论、局限和后续方向。
